\subsection{The five frontier rows in moving coordinates}
\label{sec:terminal-modal-static-frontier-reduction}

The folded source ledger above stops at $j=m-6$.  The five omitted
inequalities can be written as a six-coordinate moving-frontier system.  This
is an exact reduction; the signs of the resulting scalar expressions are not
asserted here.

Retain the rescaled correction $\bar c_n$ and inputs $E_n,F_n$ from
\eqref{eq:terminal-modal-finite-rescaling}--
\eqref{eq:terminal-modal-finite-rescaled-recurrence}.  In coordinates moving
with the newest row put
\begin{equation}
 x_n(r)=\bar c_n(n-r),\qquad
 k_n(r)=x_n(r)-\frac12x_{n-1}(r-1)\quad(2\le r\le n).
 \label{eq:terminal-modal-static-frontier-xk}
\end{equation}
Thus $k_n(r)$ is precisely the rescaled shared correction difference at
coordinate $n-r$.  Define
\begin{equation}
 z_n(0)=0,\qquad z_n(1)=2x_n(1)+E_n,\qquad
 z_n(r)=2k_n(r)\quad(r\ge2).
 \label{eq:terminal-modal-static-frontier-z}
\end{equation}
The correction recurrence becomes the exact moving-frame system
\begin{align}
 x_{n+1}(1)&=\frac12x_n(1)+F_n,
 \label{eq:terminal-modal-static-frontier-new}\\
 x_{n+1}(r)&=\frac14\{z_n(r)+2z_n(r-1)+z_n(r-2)\},
 &&2\le r\le n.
 \label{eq:terminal-modal-static-frontier-old}
\end{align}
Only the first six moving coordinates are needed below.

\begin{lemma}[Exact five-frontier reduction; proved here]
\label{lem:terminal-modal-static-five-frontier-reduction}
Put $s=\delta^{m-1}$ and, for $1\le r\le5$, let
\begin{equation}
 D_r=x_m(r)-x_{m+1}(r+1).
 \label{eq:terminal-modal-static-frontier-Dr-definition}
\end{equation}
Then
\begin{align}
 D_1&=-\frac12\{k_m(2)+E_m\},
 \label{eq:terminal-modal-static-frontier-D1}\\
 D_2&=x_m(2)-\frac12k_m(3)-k_m(2)
       -\frac12x_m(1)-\frac14E_m,
 \label{eq:terminal-modal-static-frontier-D2}\\
 D_r&=\frac12\{x_{m-1}(r-1)-k_m(r+1)-k_m(r-1)\},
 &&3\le r\le5.
 \label{eq:terminal-modal-static-frontier-Dr}
\end{align}
Moreover, if $p_m=P_m/q^2$, $S_m=\sigma_m/q^3$, and
\begin{equation}
 \beta=\eta(p_m-S_m)+\frac{\delta^m}{q\,2^m},
 \label{eq:terminal-modal-static-frontier-beta}
\end{equation}
then the five still-open rows in
\eqref{eq:terminal-modal-static-two-step-frontier}--
\eqref{eq:terminal-modal-static-two-step-cone} are exactly
\begin{align}
 \frac{d_{m-1}+d_m/4}{q^3}
   &=sD_1+\frac14\beta,
 \label{eq:terminal-modal-static-frontier-margin-1}\\
 \frac{d_{m-2}-d_m/8}{q^3}
   &=sD_2-\frac18\beta,
 \label{eq:terminal-modal-static-frontier-margin-2}\\
 \frac{d_{m-r}-d_{m-r+2}/4}{q^3}
   &=s\left(D_r-\frac14D_{r-2}\right),
 &&3\le r\le5.
 \label{eq:terminal-modal-static-frontier-margin-r}
\end{align}
Consequently, proving the nonnegativity of the five right sides above is
equivalent to closing every frontier row omitted by the folded source ledger.
No modal estimate or additional replay assumption is hidden in this
equivalence.
\end{lemma}

\begin{proof}
On an old nonseed row, the reflected $B_0$ stencil in
\eqref{eq:terminal-modal-finite-rescaled-recurrence} is the ordinary
three-point stencil.  At moving distance one its input is
$2x_n(1)+E_n$; at every larger moving distance it is
$2x_n(r)-x_{n-1}(r-1)=2k_n(r)$.  Zero extension beyond the moving frontier
gives $z_n(0)=0$.  This proves
\eqref{eq:terminal-modal-static-frontier-new}--
\eqref{eq:terminal-modal-static-frontier-old}.

Apply \eqref{eq:terminal-modal-static-frontier-old} at time $m$.
For $r=1$ it gives
\[
 x_{m+1}(2)=\frac12k_m(2)+x_m(1)+\frac12E_m,
\]
and for $r=2$ it gives
\[
 x_{m+1}(3)=\frac12k_m(3)+k_m(2)
             +\frac12x_m(1)+\frac14E_m.
\]
Subtracting these values from $x_m(1)$ and $x_m(2)$ proves
\eqref{eq:terminal-modal-static-frontier-D1}--
\eqref{eq:terminal-modal-static-frontier-D2}.  For $r\ge3$, all three
stencil entries are shared, so
\[
 x_{m+1}(r+1)=\frac12k_m(r+1)+k_m(r)+\frac12k_m(r-1).
\]
Using $x_m(r)=k_m(r)+x_{m-1}(r-1)/2$ proves
\eqref{eq:terminal-modal-static-frontier-Dr}.

The literal final degree change affects only the new endpoint correction.
Hence, on every old row $j=m-r$,
\[
 d_j=\delta C_m(j)-C_{m+1}^{\rm full}(j)
     =q^3\delta^{m-1}D_r.
\]
At the endpoint, \eqref{eq:terminal-modal-static-b} and
\eqref{eq:terminal-modal-directed-deconvolution} give
\[
 \frac{d_m}{q^3}
 =\eta(p_m-S_m)+\frac{\delta^m}{q\,2^m}=\beta.
\]
Substitution in the five two-step inequalities proves
\eqref{eq:terminal-modal-static-frontier-margin-1}--
\eqref{eq:terminal-modal-static-frontier-margin-r}.
\end{proof}

The focused exact preflight checks the moving-frame recurrence and every
identity in \cref{lem:terminal-modal-static-five-frontier-reduction} against
literal rational replays.  Finite screens show that the second expression
\eqref{eq:terminal-modal-static-frontier-margin-2} is the tightest of the
five.  All five are only of order $q$ after division by $q^3$, whereas the
individual terms in their formulas are of constant order.  Thus the proved
sign $k_n(r)>0$ alone cannot close them: the remaining task is a uniform
first-order comparison of the coupled $E_n,F_n$ frontier history and the
endpoint scalar $S_m$.  The exact reduction and the finite screen do not
promote that comparison.

\subsubsection{The joint first-order limit}

The cancellation just described is exact at leading order.  The next
coefficient nevertheless has a simple positive form.  Besides proving that
the five rows are eventually positive, the calculation below isolates an
exact finite remainder interface for a future uniform estimate.

For $r\ge1$ put
\begin{equation}
 c_r=-\frac43+\frac13\left(-\frac12\right)^r,
 \qquad c_0=0.
 \label{eq:terminal-modal-static-frontier-cr}
\end{equation}
For every index at which the quantities are defined, set
\begin{equation}
 a_n=\frac{E_{n+1}-E_n}{q},\qquad
 \nu_n=-\frac{\mu_n}{q},\qquad
 y_n(r)=\frac{x_n(r)-c_rE_n}{q}.
 \label{eq:terminal-modal-static-frontier-finite-interface}
\end{equation}
Thus no asymptotic quantity occurs in the definition of $y_n$.  The exact
moving recurrence gives
\begin{align}
 y_{n+1}(1)&=\frac12y_n(1)-\nu_n-c_1a_n,
 \label{eq:terminal-modal-static-frontier-y1}\\
 y_{n+1}(2)&=\frac14\{2y_n(2)+4y_n(1)-y_{n-1}(1)
              +c_1a_{n-1}\}-c_2a_n,
 \label{eq:terminal-modal-static-frontier-y2}\\
 y_{n+1}(r)&=\frac14\bigl\{
 2y_n(r)+4y_n(r-1)+2y_n(r-2)
 -y_{n-1}(r-1)-2y_{n-1}(r-2)-y_{n-1}(r-3)\notag\\
 &\hspace{35mm}
 +(c_{r-1}+2c_{r-2}+c_{r-3})a_{n-1}\bigr\}-c_ra_n,
 \qquad r\ge3.
 \label{eq:terminal-modal-static-frontier-yr}
\end{align}
Here $y_n(0)=0$.  These equations are an exact finite triangular system, not
a formal expansion.

Define
\begin{equation}
 \ell_{m,r}=y_m(r)-y_{m+1}(r+1)-c_{r+1}a_m,
 \qquad
 \mathcal B_m=\frac{\beta-\delta^{m-1}E_m}{q}.
 \label{eq:terminal-modal-static-frontier-ellB}
\end{equation}
If $\mathcal M_{m,1},\ldots,\mathcal M_{m,5}$ denote, in order, the five
left sides before division by $q^3$ in
\eqref{eq:terminal-modal-static-frontier-margin-1}--
\eqref{eq:terminal-modal-static-frontier-margin-r}, then the cancellation of
$c_r-c_{r+1}=\frac12(-1/2)^r$ gives the exact identities
\begin{align}
 \frac{\mathcal M_{m,1}}{q^4}
 &=\delta^{m-1}\ell_{m,1}+\frac14\mathcal B_m,
 \label{eq:terminal-modal-static-frontier-finite-M1}\\
 \frac{\mathcal M_{m,2}}{q^4}
 &=\delta^{m-1}\ell_{m,2}-\frac18\mathcal B_m,
 \label{eq:terminal-modal-static-frontier-finite-M2}\\
 \frac{\mathcal M_{m,r}}{q^4}
 &=\delta^{m-1}\left(\ell_{m,r}-\frac14\ell_{m,r-2}\right),
 \qquad 3\le r\le5.
 \label{eq:terminal-modal-static-frontier-finite-Mr}
\end{align}
Equations \eqref{eq:terminal-modal-static-frontier-y1}--
\eqref{eq:terminal-modal-static-frontier-finite-Mr} are the promised uniform
remainder interface: it is enough to control the local slopes $a_n$, masses
$\nu_n$, the contractive triangular transient, and the single endpoint
quantity $\mathcal B_m$.

\begin{theorem}[Positive joint first-order frontier limit; proved here]
\label{thm:terminal-modal-static-frontier-first-order}
Along $q=1/(16m)$, let $m\to\infty$ and put $T=e^{-1/8}$.  Define
\begin{align}
 f(t)&=\frac{t^4-30t^3+12t^2+10t+3}
 {10(1+t^2)^2},\notag\\
 g(t)&=\frac{t^4-10t^3+6t^2+10t+1}
 {5(1+t^2)^2}.
 \label{eq:terminal-modal-static-frontier-fg}
\end{align}
Then
\begin{align}
 \lim_{m\to\infty}\frac{\mathcal M_{m,1}}{q^4}
 &=2\left(g(T)-\frac{13}{8}f(T)\right)>\frac23,
 \label{eq:terminal-modal-static-frontier-limit-1}\\
 \lim_{m\to\infty}\frac{\mathcal M_{m,2}}{q^4}
 &=g(T)-\frac74f(T)>\frac13,
 \label{eq:terminal-modal-static-frontier-limit-2}\\
 \lim_{m\to\infty}\frac{\mathcal M_{m,3}}{q^4}
 &=g(T)-\frac{13}{8}f(T)>\frac13,
 \label{eq:terminal-modal-static-frontier-limit-3}\\
 \lim_{m\to\infty}\frac{\mathcal M_{m,4}}{q^4}
 &=g(T)-\frac{27}{16}f(T)>\frac13,
 \label{eq:terminal-modal-static-frontier-limit-4}\\
 \lim_{m\to\infty}\frac{\mathcal M_{m,5}}{q^4}
 &=g(T)-\frac{53}{32}f(T)>\frac13.
 \label{eq:terminal-modal-static-frontier-limit-5}
\end{align}
In particular, all five frontier rows are strictly positive for all
sufficiently large $m$.  This theorem does not identify an explicit cutoff;
the desired uniform statement beginning at $m=64$ remains open.
\end{theorem}

\begin{proof}
Write $\tau=nq$ and $t=e^{-2\tau}$.  From
\eqref{eq:terminal-modal-static-frontier-function} and
\eqref{eq:terminal-modal-finite-E}, uniformly for a fixed number of indices
about $n=m$,
\begin{align}
 p_n&\longrightarrow
 p(\tau)=\frac{2(t^2+10t-1)}{5(1+t^2)},\notag\\
 E_n&\longrightarrow E(\tau)=e^\tau e_0(t),
 \qquad
 e_0(t)=\frac{t^2-10t+3}{10(1+t^2)}.
 \label{eq:terminal-modal-static-frontier-input-limits}
\end{align}
Consequently $a_n\to E'(\tau)=e^\tau f(t)$.  Formula
\eqref{eq:terminal-modal-finite-mu-M}, together with
\eqref{eq:terminal-modal-finite-mass-ratio}, gives
\begin{equation}
 \nu_n\longrightarrow e^\tau g(t).
 \label{eq:terminal-modal-static-frontier-nu-limit}
\end{equation}
Indeed, the undamped mass bracket tends to
\[
 \frac{p(\tau)}4\frac{2(1-t^2)}{1+t^2}+\frac25,
\]
which simplifies to $g(t)$.  Direct differentiation of $e_0$ gives
\[
 e_0(t)-2te_0'(t)=f(t).
\]

The coefficient $1/2$ in
\eqref{eq:terminal-modal-static-frontier-y1}, followed by the triangular
structure of \eqref{eq:terminal-modal-static-frontier-y2}--
\eqref{eq:terminal-modal-static-frontier-yr}, makes every initial transient
vanish as $n\to\infty$.  More explicitly, at each fixed moving distance $r$
the homogeneous transient is bounded by a fixed polynomial in $n$ times
$2^{-n}$; the initial values grow at most polynomially in $1/q$, while the
recent forcing converges to the displayed local limits.  Substitution of the limits
$\nu=e^\tau g(t)$ and $a=e^\tau f(t)$ gives, for $1\le r\le7$,
\begin{equation}
\begin{array}{c|rrrrrrr}
 r&1&2&3&4&5&6&7\\ \hline
 \operatorname{coeff}_{\nu}y(r)&-2&-3&-9/2&-23/4&-57/8&-135/16&-313/32\\
 \operatorname{coeff}_{a}y(r)&3&25/4&10&215/16&273/16&1317/64&773/32.
\end{array}
 \label{eq:terminal-modal-static-frontier-y-limit-table}
\end{equation}
The table follows successively from the exact triangular recurrence; the
focused verifier performs the same calculation over the rationals.  Hence
\begin{equation}
 \ell_{m,r}\longrightarrow L_r,\qquad
\begin{array}{c|rrrrr}
 r&1&2&3&4&5\\ \hline
 \operatorname{coeff}_{\nu}L_r&1&3/2&5/4&11/8&21/16\\
 \operatorname{coeff}_{a}L_r&-2&-19/8&-17/8&-73/32&-35/16.
\end{array}
 \label{eq:terminal-modal-static-frontier-L-limit-table}
\end{equation}

It remains to identify the endpoint term.  Put
$\zeta_n=S_n-(3p_n/2-2/5)$.  The exact recurrence
\eqref{eq:terminal-modal-static-endpoint-error} and its contraction factor
$\eta/(1+q)=\delta/2$ imply
\begin{equation}
 \frac{\zeta_m}{q}\longrightarrow
 -p'(\tau)-3p(\tau)+\frac45.
 \label{eq:terminal-modal-static-frontier-zeta-limit}
\end{equation}
The initial term is exponentially small, while the forcing divided by $q$
tends to $-p'/2-3p/2+2/5$; the factor $(1-1/2)^{-1}$ supplies the displayed
limit.  There is also the exact simplification
\begin{equation}
 \mathcal B_m=
 \frac{\eta}{2q}\{\delta p_{m-1}-p_m\}+\frac\delta5
 -\frac\eta q\zeta_m+\frac{\delta^m}{q^2 2^m}.
 \label{eq:terminal-modal-static-frontier-B-exact}
\end{equation}
It follows that
\begin{equation}
 \mathcal B_m\longrightarrow4g(t)-5f(t).
 \label{eq:terminal-modal-static-frontier-B-limit}
\end{equation}
For the last simplification use
$p=4/5-4e_0$, $p'=4(e_0-f)$, and
\begin{equation}
 g=f+\frac15-e_0.
 \label{eq:terminal-modal-static-frontier-g-identity}
\end{equation}

At $\tau=1/16$, $\delta^{m-1}\to e^{-1/16}$.  Thus the exponential in
\eqref{eq:terminal-modal-static-frontier-nu-limit}--
\eqref{eq:terminal-modal-static-frontier-L-limit-table} cancels.  Substituting
the two tables and \eqref{eq:terminal-modal-static-frontier-B-limit} into
\eqref{eq:terminal-modal-static-frontier-finite-M1}--
\eqref{eq:terminal-modal-static-frontier-finite-Mr} proves the five stated
limits.

Finally $7/8<T<1$.  Let
\[
 F=t^4-30t^3+12t^2+10t+3,\qquad
 G=t^4-10t^3+6t^2+10t+1
\]
and put $H_2=8G-7F$, $H_3=16G-13F$,
$H_4=32G-27F$, and $H_5=64G-53F$.  With
$y=8t-7\in[0,1]$, direct substitution gives
\begin{align*}
 4096H_2((7+y)/8)
 =228817+127116y+19830y^2+1068y^3+y^4,\\
 4096H_3((7+y)/8)
 =463475+236196y+35682y^2+1924y^3+3y^4,\\
 4096H_4((7+y)/8)
 =921109+490428y+75342y^2+4060y^3+5y^4,\\
 4096H_5((7+y)/8)
 =1848059+962820y+146706y^2+7908y^3+11y^4.
\end{align*}
Every displayed coefficient is positive.  The four limits other than the
factor two in \eqref{eq:terminal-modal-static-frontier-limit-1} are,
respectively,
\[
 \frac{H_2}{40(1+t^2)^2},\qquad
 \frac{H_3}{80(1+t^2)^2},\qquad
 \frac{H_4}{160(1+t^2)^2},\qquad
 \frac{H_5}{320(1+t^2)^2}.
\]
Using $(1+t^2)^2\le4$ and the displayed constant terms proves the rational
lower bounds $1/3$ (and twice that bound in
\eqref{eq:terminal-modal-static-frontier-limit-1}).
\end{proof}
